\begin{chapterpage}{Inferensi untuk data numerik}
  \chaptertitle{Inferensi untuk data numerik}
  \label{inferenceForNumericalData}
  \label{ch_inference_for_means}
  \chaptersection{oneSampleMeansWithTDistribution}
% Boundary-clean reader omits navigation to excluded later Chapter 7 sections.
\end{chapterpage}
\renewcommand{\chapterfolder}{ch_inference_for_means}


\chapterintro{Bab~\ref{ch_foundations_for_inf}
  memperkenalkan kerangka inferensi statistik yang didasarkan
  pada interval kepercayaan dan hipotesis menggunakan
  distribusi normal untuk proporsi sampel.
  Dalam bab ini, kita akan menjumpai beberapa estimasi titik baru
  dan dua distribusi baru.
  Dalam setiap kasus, gagasan inferensinya tetap sama:
  tentukan estimasi titik atau statistik uji yang berguna,
  identifikasi distribusi yang sesuai untuk estimasi titik
  atau statistik uji tersebut, lalu
  terapkan gagasan inferensi.}



%__________________
\section[Rata-rata satu sampel dengan distribusi $t$]
    {Rata-rata satu sampel dengan distribusi
        $\pmb{\MakeLowercase{t}}$}
\label{oneSampleMeansWithTDistribution}

\noindent%
Sebagaimana kita dapat memodelkan perilaku
proporsi sampel $\hat{p}$ menggunakan distribusi normal,
rata-rata sampel $\bar{x}$ juga dapat dimodelkan menggunakan
distribusi normal ketika syarat-syarat tertentu terpenuhi.
\index{estimasi titik!rata-rata tunggal}
Namun, kita akan segera mempelajari bahwa distribusi baru,
yang disebut distribusi $t$,
cenderung lebih berguna ketika bekerja dengan rata-rata sampel.
Pertama-tama kita akan mempelajari distribusi baru ini,
kemudian menggunakannya untuk menyusun interval kepercayaan
dan melakukan uji hipotesis untuk rata-rata.


\subsection[Distribusi $\bar{x}$]
    {Distribusi sampling dari $\pmb{\bar{x}}$}

Rata-rata sampel cenderung mengikuti
distribusi normal yang berpusat pada rata-rata populasi,~$\mu$,
ketika syarat-syarat tertentu terpenuhi.
Selain itu, kita dapat menghitung galat baku untuk rata-rata
sampel menggunakan simpangan baku populasi $\sigma$
dan ukuran sampel $n$.

\begin{onebox}{Teorema Limit Pusat untuk rata-rata sampel}
  Ketika kita mengumpulkan sampel cukup besar yang terdiri atas
  $n$~observasi yang saling independen dari populasi dengan
  rata-rata $\mu$ dan simpangan baku $\sigma$,
  distribusi sampling $\bar{x}$ akan mendekati
  distribusi normal dengan
  \begin{align*}
  &\text{Rata-rata}=\mu
  &&\text{Galat Baku }(SE) = \frac{\sigma}{\sqrt{n}}
  \end{align*}
\end{onebox}

\noindent%
Sebelum mendalami interval kepercayaan dan uji
hipotesis menggunakan $\bar{x}$, pertama-tama kita perlu membahas dua topik:
\begin{itemize}
\item
    Ketika kita memodelkan $\hat{p}$ menggunakan distribusi normal,
    syarat-syarat tertentu harus terpenuhi.
    Syarat-syarat untuk bekerja dengan $\bar{x}$
    sedikit lebih rumit, dan kita akan menggunakan
    Bagian~\ref{x_bar_conditions} untuk membahas
    cara memeriksa syarat-syarat inferensi.
\item
    Galat baku bergantung pada simpangan
    baku populasi, $\sigma$.
    Namun, kita jarang mengetahui $\sigma$, sehingga
    kita harus mengestimasinya.
    Karena estimasi ini sendiri tidak sempurna,
    kita menggunakan distribusi baru yang disebut
    distribusi $t$\index{distribusi t@distribusi $t$}
    untuk mengatasi masalah ini, yang kita bahas dalam
%    While we can use the plug-in principle,
%    using the sample standard deviation $s$ in place of $\sigma$,
%    this is not quite enough to resolve the issue entirely.
%    and .
%    We'll cover this topic in
    Bagian~\ref{introducingTheTDistribution}.
\end{itemize}


\subsection[Mengevaluasi dua syarat yang diperlukan untuk
    memodelkan $\bar{x}$]
  {Mengevaluasi dua syarat yang diperlukan untuk
      memodelkan $\pmb{\bar{x}}$}
\label{x_bar_conditions}

\noindent%
Dua syarat diperlukan untuk menerapkan
Teorema Limit Pusat\index{Teorema Limit Pusat}
pada rata-rata sampel~$\bar{x}$:
\begin{description}
\item[Independensi.]
    Observasi sampel harus saling independen.
    Cara paling umum untuk memenuhi syarat ini adalah
    mengambil sampel acak sederhana dari
    populasi.
    Jika data berasal dari proses acak,
    yang serupa dengan pelemparan sebuah dadu,
    hal ini juga akan memenuhi syarat independensi.
\item[Normalitas.]
    Ketika sampelnya kecil,
    kita juga mensyaratkan bahwa observasi sampel
    berasal dari populasi yang berdistribusi normal.
    Syarat ini semakin dapat dilonggarkan
    seiring bertambahnya ukuran sampel.
    Syarat ini memang bersifat samar,
    sehingga sulit dievaluasi;
    karena itu, selanjutnya kita memperkenalkan dua pedoman praktis
    agar pemeriksaan syarat ini menjadi lebih mudah.
\end{description}
%%Before we get to the sample size consideration, let's
%%consider a special case of the normal distribution
%%where any sample size is sufficient.
%
%%There is also a special case of the Central Limit Theorem
%%for when the data come from a nearly normal distribution.
%%In this case the sample mean will be nearly normal
%%regardless of sample size.
%
%\begin{onebox}{Special case of the Central Limit Theorem
%    for normally distributed data}
%  The sampling distribution of $\bar{x}$ is nearly normal when
%  the sample observations are independent and come from a nearly
%  normal distribution.
%  This is true for any sample size.
%\end{onebox}
%
%%For population distributions that are not normal,
%%the sample mean $\bar{x}$ will still look normal if the sample
%%size is large enough.
%%To check what is \emph{large enough}, we ask two questions:
%%\begin{itemize}
%%\item
%%    Is the sample show evident skew or outliers?
%%    If so, then if t
%%\end{itemize}
%
%In practice, the population never exactly follows
%a normal distribution,
%and the more ``non-normal'' a population
%distribution, the larger the required sample size required for
%$\bar{x}$ to be reasonably modeled using a normal distribution.
%The rough rule of thumb is, if you don't see any clear outliers
%and we don't have reason to believe particularly extreme outliers
%are present in population, then this condition is satisfied.

\begin{onebox}{Pedoman praktis:
    cara melakukan pemeriksaan normalitas}
  Tidak ada cara sempurna untuk memeriksa syarat normalitas,
  jadi sebagai gantinya kita menggunakan dua pedoman praktis: %,
%  one for small samples ($n < 30$)
%  and another for large samples ($n \geq 30$):
  \begin{description}
  \setlength{\itemsep}{0mm}
  \item[$\mathbf{n < 30}$:]
      Jika ukuran sampel $n$ kurang dari 30
      dan tidak terdapat pencilan yang jelas dalam data,
      biasanya kita mengasumsikan data berasal dari
      distribusi yang mendekati normal agar
      syarat tersebut terpenuhi.
  \item[$\mathbf{n \geq 30}$:]
      Jika ukuran sampel $n$ setidaknya 30
      dan tidak terdapat pencilan yang \emph{sangat ekstrem},
      biasanya kita mengasumsikan distribusi sampling
      $\bar{x}$ mendekati normal, sekalipun distribusi
      observasi individual yang mendasarinya tidak normal.
  \end{description}
\end{onebox}

Dalam mata kuliah statistika pengantar ini, Anda tidak diharapkan
mengembangkan penilaian yang sempurna mengenai syarat normalitas.
Namun, Anda diharapkan mampu menangani
kasus-kasus yang jelas berdasarkan pedoman praktis tersebut.\footnote{Pedoman
  yang lebih cermat akan mempertimbangkan pelonggaran lebih lanjut
  atas pemeriksaan \emph{pencilan sangat ekstrem} ketika
  ukuran sampel sangat besar.
  Namun, pembahasan lebih lanjut akan kita sisakan untuk mata kuliah berikutnya.}

\begin{examplewrap}
\begin{nexample}{Perhatikan dua grafik berikut
    yang berasal dari sampel acak sederhana dari
    populasi yang berbeda.
    Ukuran sampelnya adalah $n_1 = 15$ dan $n_2 = 50$.
    \begin{center}
    \Figure[Dua histogram ditampilkan, satu untuk "Observasi Sampel 1" dan satu lagi untuk "Observasi Sampel 2". Histogram Observasi Sampel 1 memuat nilai dari 0 hingga 7 dengan lebar bin 1, sehingga terdapat 7 bin dengan frekuensi 2, 1, 4, 3, 2, 0, dan 3. Histogram Observasi Sampel 2 memuat nilai dari 0 hingga 22, dengan lebar bin 1. Sebagian besar data berada dekat nol, dengan separuh observasi berada dalam bin dari 0 hingga 1. Hanya ada satu bin tak nol di atas 5; tingginya tampak 1, yaitu bin dari 21 hingga 22.]{0.85}{outliers_and_ss_condition}
    \end{center}
    Apakah syarat independensi dan normalitas terpenuhi
    dalam setiap kasus?}
  \label{outliers_and_ss_condition_ex}%
  Setiap sampel merupakan sampel acak sederhana dari
  populasinya masing-masing, sehingga syarat independensi
  terpenuhi.
  Selanjutnya, mari kita periksa syarat normalitas untuk
  setiap sampel menggunakan pedoman praktis.
  
  Sampel pertama memiliki kurang dari 30 observasi,
  sehingga kita memperhatikan setiap pencilan yang jelas.
  Tidak ada pencilan semacam itu; meskipun terdapat celah kecil dalam
  histogram antara 5 dan~6, celah ini kecil dan
  20\% observasi dalam sampel kecil ini
  terwakili oleh batang paling kanan pada histogram,
  sehingga observasi itu nyaris tidak dapat disebut pencilan yang jelas.
  Karena tidak ada pencilan yang jelas, syarat normalitas
  secara wajar~terpenuhi.

  Sampel kedua berukuran lebih besar dari 30 dan
  mencakup pencilan yang tampaknya kira-kira 5 kali
  lebih jauh dari pusat distribusi dibandingkan
  observasi terjauh berikutnya.
  Ini merupakan contoh pencilan yang sangat ekstrem,
  sehingga syarat normalitas tidak akan terpenuhi.
\end{nexample}
\end{examplewrap}

Dalam praktiknya, lazim pula melakukan pemeriksaan secara mental untuk menilai
apakah ada alasan untuk meyakini populasi yang mendasari
memiliki kemencengan sedang (jika $n < 30$)
atau memiliki pencilan yang sangat ekstrem ($n \geq 30$)
di luar apa yang kita amati dalam data.
Sebagai contoh, perhatikan jumlah pengikut
pada setiap akun individual di Twitter,
lalu bayangkan distribusinya.
Sebagian besar akun telah mengumpulkan
paling banyak beberapa ribu pengikut,
sementara sebagian yang relatif sangat kecil telah menghimpun
puluhan juta pengikut,
yang berarti distribusinya sangat menceng.
Ketika kita mengetahui bahwa data berasal dari distribusi
yang kemencengannya seekstrem itu,
diperlukan upaya tertentu untuk menentukan seberapa besar
ukuran sampel yang diperlukan agar syarat normalitas
terpenuhi.

%if we were sampling accounts from Twitter
%and examining the distribution of followers on the sampled
%accounts, we can expect that the vast majority of accounts
%will have fewer than 1,000 followers and that there
%will be some very extreme outliers who have tens of millions
%of followers.

%Distribution of the number of subscribers for
%      anyone who has uploaded a video to YouTube.
%      Most such individuals will have built little to
%      no following, while others will have amassed tens
%      of millions of subscribers.
%Generally, we do not presume you to always know when the
%underlying population has particularly extreme outliers.
%That~is, besides looking at the data itself,
%considering the mental check for whether particularly extreme
%outliers are likely to be a sanity check, not a formal check.

%\begin{figure}[h]
%  \centering
%  \Figure{0.8}{outliers_and_ss_condition}
%  \caption{Sample observations for
%      Example~\ref{outliers_and_ss_condition_ex}.}
%  \label{outliers_and_ss_condition}
%\end{figure}

%A more thorough sample size condition assessment would
%also consider two additional aspects beyond the core
%guidance above.

%The most nuanced checks are then when the sample size
%is very small -- and we have almost no observational data
%to allow us to check the condition.
%\begin{description}
%\item[Population knowledge.]
%    If we have information about the population beyond
%    what we've observed in the sample, we would consider
%    this information as well.
%    For example, if the sample size is under 30
%    and the population is known to be moderately skewed
%    (something difficult to detect with a small sample
%    in the observed data),
%    we might still not consider the 
%    For example, if the sample size is under 30 but the
%    population is known to be moderately skewed,
%    then the sample size condition is not reasonable.
%    Likewise, if the population is known to have particularly
%    extreme observations (examples below), then we may
%    require a particularly large sample size if we want
%    to use the normal model for $\bar{x}$.
%\item[Relaxing the extreme outlier condition.]
%    When the sample size gets very large,
%    we may even be able to overcome issues with
%    particularly extreme outliers.
%    However, there isn't clear guidance, and instead,
%    custom simulations can be helpful but are beyond
%    the scope of this book.
%\end{description}
%In this first course in statistics,
%you won't (and aren't expected to) have perfect judgement
%on when the sample size condition is or is not met.
%However, you are expected to be able to handle the
%clear cut cases based on the core guidelines.

%For those wanting to do more rigorous checks
%or for the situation that the ,
%then we add a 
%below are slightly more careful checks, it's convenient
%to break down 
%\begin{description}
%\item[Sample size under 30.]
%    If the sample size is less than 30, then we simply follow
%    the rule of thumb and there isn't
%    extreme skew in the data (usually punctuated by
%    extreme outliers), then we can proceed.
%\item[Sample size at least 30.]
%    If the sample size is at least 30 and there isn't
%    extreme skew in the data (usually punctuated by
%    extreme outliers), then we can proceed.
%\end{description}
%then it's generally reasonable to consider $\bar{x}$
%as following a nearly normal distribution.

%\Comment{Check the ``99\%'' and ``hundreds'' claim in the
%  income example below.}
%
%\begin{examplewrap}
%\begin{nexample}{Describe a couple populations that you know
%    would have particularly extreme outliers.}
%  Wealth distributions in many countries have
%      particularly extreme outliers.
%      For example, over 99\% of the population
%      has fewer than \$10 million saved,
%      while there are hundreds individuals in the
%      United States with over \$1~billion and who
%      are unlikely to be captured in even a moderate-sized
%      sample.
%
%  Distribution of the number of subscribers for
%      anyone who has uploaded a video to YouTube.
%      Most such individuals will have built little to
%      no following, while others will have amassed tens
%      of millions of subscribers.
%
%%  So while we won't be quizzing you on a variety of applications
%%  in this book, when you apply these skills elsewhere it is
%%  important to keep this consideration in mind and do some
%%  research if you aren't sure about outliers.
%\end{nexample}
%\end{examplewrap}
%
%Generally, we do not presume you to always know when the
%underlying population has particularly extreme outliers.
%That~is, besides looking at the data itself,
%considering the mental check for whether particularly extreme
%outliers are likely to be a sanity check, not a formal check.

%\begin{examplewrap}
%\begin{nexample}{Suppose we randomly sampled 20 individuals
%    from the United States and considered their incomes. %,
%    % which are shown in the following distribution:
%    However, the population is known to have particularly
%    extreme outliers, e.g. some individuals with incomes
%    above \$10 million.
%    No matter what we observe in the original 20 observations,
%    can you say whether we should proceed with modeling
%    $\bar{x}$ using a normal distribution?}
%  When we know the population distribution to have particularly
%  extreme outliers, then even if we observe no outliers in our
%  sample, we should not proceed to model $\bar{x}$ using
%  a normal distribution.
%
%  Generally, we do not presume you to always know when the
%  underlying population has particularly extreme outliers.
%  So while we won't be quizzing you on a variety of applications
%  in this book, when you apply these skills elsewhere it is
%  important to keep this consideration in mind and do some
%  research if you aren't sure about outliers.
%\end{nexample}
%\end{examplewrap}

%However, if one or more of clear outliers are present are evidently present,
%the guidelines around a reasonable minimum sample become murky.
%If 
%We'll see some other examples throughout the rest of this book,
%which will help in developing some intuition around this topic,
%but in many cases, data with .

\index{Teorema Limit Pusat!data normal|)}



\subsection[Memperkenalkan distribusi $t$]
    {Memperkenalkan distribusi $\pmb{t}$}
\label{introducingTheTDistribution}

\index{distribusi t@distribusi $t$|(}
\index{distribusi!t@$t$|(}

Dalam praktiknya, kita tidak dapat menghitung galat baku secara langsung
untuk $\bar{x}$ karena kita tidak mengetahui simpangan baku
populasi,~$\sigma$.
Kita menghadapi masalah serupa ketika menghitung galat
baku untuk proporsi sampel, yang bergantung pada proporsi
populasi,~$p$.
Solusi kita dalam konteks proporsi adalah menggunakan suatu nilai
sampel sebagai pengganti
nilai populasi ketika menghitung galat baku.
Kita akan memakai strategi serupa untuk menghitung galat
baku $\bar{x}$, dengan menggunakan simpangan
baku sampel $s$ sebagai pengganti $\sigma$:
\begin{align*}
SE = \frac{\sigma}{\sqrt{n}} \approx \frac{s}{\sqrt{n}}
\end{align*}
Strategi ini cenderung bekerja baik ketika kita memiliki
banyak data dan dapat mengestimasi $\sigma$ secara akurat menggunakan $s$.
Namun, estimasi tersebut kurang presisi pada sampel yang lebih kecil,
dan hal ini menimbulkan masalah ketika distribusi
normal digunakan untuk memodelkan $\bar{x}$.
% --
%when the sample size is large --
%but it is less reliable when the sample size is smaller
%than about 30. % independent observations.

Kita akan mendapati bahwa distribusi baru yang disebut
\termsub{distribusi $\pmb{t}$}{distribusi t@distribusi $t$} berguna
untuk perhitungan inferensi.
Distribusi~$t$, yang ditampilkan sebagai garis penuh dalam
Gambar~\ref{tDistCompareToNormalDist}, berbentuk lonceng.
Namun, ekornya lebih tebal daripada ekor distribusi normal,
yang berarti observasi lebih mungkin jatuh lebih dari dua
simpangan baku dari rata-rata dibandingkan pada distribusi
normal. %\footnote{The standard deviation of the
  %$t$-distribution is actually a little more than 1.
  %However, it is useful to always think of the $t$-distribution
  %as having a standard deviation of 1 in all of our applications.}
%This distribution is important since it accounts for
%a key challenge with modeling the sample mean:
%the standard error of the sample mean isn't as
%precise when the sample size is small.
Ketebalan ekstra pada ekor distribusi $t$ merupakan
koreksi yang diperlukan untuk mengatasi masalah penggunaan~$s$
sebagai pengganti $\sigma$ dalam perhitungan $SE$.

\begin{figure}[h]
  \centering
  \Figure[Ditampilkan distribusi normal baku dan distribusi t. Distribusi t juga berbentuk lonceng, tetapi puncaknya lebih rendah dan lebih landai daripada distribusi normal serta ekornya lebih tebal. Sebagai contoh, terdapat bagian distribusi yang cukup besar -- mungkin 5\% untuk distribusi t khusus ini -- yang membentang hingga di bawah -3 dan di atas 3, sedangkan distribusi normal sangat mendekati nol pada nilai di bawah -3 atau di atas 3.]{0.7}{tDistCompareToNormalDist}
  \caption{Perbandingan distribusi $t$
      dan distribusi normal.}
  \label{tDistCompareToNormalDist}
\end{figure}

Distribusi $t$ selalu berpusat di nol dan
memiliki satu parameter: derajat kebebasan.
\termsub{Derajat kebebasan ($\pmb{df}$)}
    {derajat kebebasan ($df$)!distribusi $t$}
menentukan bentuk tepat distribusi $t$ yang berbentuk lonceng.
Beberapa distribusi $t$ ditampilkan dalam
Gambar~\ref{tDistConvergeToNormalDist}
sebagai perbandingan dengan distribusi normal.

Secara umum, kita akan menggunakan distribusi $t$
dengan $df = n - 1$ untuk memodelkan rata-rata sampel
ketika ukuran sampelnya $n$.
Artinya, ketika kita memiliki lebih banyak observasi,
derajat kebebasannya akan lebih besar dan
distribusi $t$ akan semakin menyerupai
distribusi normal baku;
ketika banyaknya derajat kebebasan sekitar 30 atau lebih,
distribusi $t$ hampir tidak dapat dibedakan
dari distribusi normal.

\begin{figure}[h]
  \centering
  \Figure[Empat distribusi t dengan derajat kebebasan 1, 2, 4, dan 8 ditampilkan bersama distribusi normal dalam grafik yang sama. Semakin besar derajat kebebasan, semakin dekat distribusi t dengan distribusi normal; puncak di pusat semakin tinggi dan mendekati puncak distribusi normal, sedangkan ekor distribusinya tampak semakin tipis.]{0.75}{tDistConvergeToNormalDist}
  \caption{Semakin besar derajat kebebasan,
      distribusi $t$ semakin menyerupai distribusi
      normal baku.}
  \label{tDistConvergeToNormalDist}
\end{figure}

\begin{onebox}{Derajat kebebasan
    ($\pmb{\MakeLowercase{df}}$)}
  Derajat kebebasan menentukan bentuk
  distribusi $t$.
  Semakin besar derajat kebebasan, semakin baik
  distribusi tersebut menghampiri model normal. \stdvspace{}

  Ketika memodelkan $\bar{x}$ menggunakan distribusi $t$,
  gunakan $df = n - 1$.
\end{onebox}

%\Comment{Cut this next sentence?}
%In Section~\ref{tDistSolutionToSEProblem},
%we relate degrees of freedom to sample size.

Distribusi $t$ memberi kita fleksibilitas lebih besar daripada
distribusi normal ketika menganalisis data numerik.
Dalam~praktiknya, perangkat lunak statistik lazim digunakan,
seperti R, Python, atau SAS untuk analisis ini.
Sebagai alternatif, kalkulator grafik atau
\termsub{tabel $\pmb{t}$}{tabel t@tabel $t$} dapat digunakan;
tabel $t$ serupa dengan tabel distribusi normal,
dan dapat ditemukan dalam Lampiran~\ref{tDistributionTable},
yang memuat petunjuk penggunaan dan contoh
bagi mereka yang ingin menggunakan opsi ini.
Apa pun pendekatan yang dipilih, terapkan metode Anda
pada contoh-contoh di bawah ini untuk memastikan
pemahaman Anda mengenai distribusi $t$.

\begin{examplewrap}
\begin{nexample}{Berapa proporsi distribusi $t$
    dengan 18 derajat kebebasan yang berada di bawah -2.10?}
  Seperti dalam soal peluang normal, pertama-tama kita menggambar
  grafik dalam Gambar~\ref{tDistDF18LeftTail2Point10}
  dan mengarsir daerah di bawah -2.10.
%  If this were a normal distribution, the area would be
%  a little less than 0.025, since about 5\% of the area
%  under a normal curve goes out beyond $\pm 1.96$ standard
%  deviations.
  Dengan perangkat lunak statistik, kita dapat memperoleh
  nilai yang presisi: 0.0250.
%  The tail area below -2.10 in the $t$-distribution with
%  $df = 18$ is the same as the tail area below -1.96 in
%  the normal distribution.
\end{nexample}
\end{examplewrap}

\begin{figure}
  \centering
  \Figure[Ditampilkan distribusi t dengan 18 derajat kebebasan; daerah di bawah -2.10 diarsir dan tampak mewakili sekitar 2\% hingga 5\% dari luas distribusi. Secara umum, ketika derajat kebebasan lebih besar daripada sekitar 10, seperti dalam kasus ini, perbedaan visual antara distribusi t dan distribusi normal sangat halus, meskipun perbedaan tersebut tetap penting untuk perhitungan kita.]{0.42}{tDistDF18LeftTail2Point10}
  \caption{Distribusi $t$ dengan 18 derajat kebebasan.
      Daerah di bawah -2.10 telah diarsir.}
  \label{tDistDF18LeftTail2Point10}
\end{figure}

\begin{examplewrap}
\begin{nexample}{Distribusi $t$ dengan 20 derajat kebebasan
    ditampilkan pada panel kiri
    Gambar~\ref{tDistDF20RightTail1Point65}.
    Estimasikan proporsi distribusi yang berada
    di atas 1.65.}
  Pada distribusi normal, nilainya akan sekitar
  sekitar~0.05, sehingga kita mengharapkan distribusi $t$
  memberikan nilai di sekitar itu.
  Dengan perangkat lunak statistik: 0.0573.
\end{nexample}
\end{examplewrap}

\begin{figure}
  \centering
  \Figure[Dua distribusi t ditampilkan pada dua grafik terpisah. Grafik pertama menampilkan distribusi t dengan 20 derajat kebebasan dan daerah di atas 1.65 diarsir; daerah itu tampak mewakili sekitar 5\% dari luas total distribusi. Grafik kedua menampilkan distribusi t dengan 2 derajat kebebasan dan daerah di bawah -3 serta di atas 3 diarsir. Karena derajat kebebasannya sangat kecil, ekor distribusi ini jauh lebih tebal dan pusatnya juga berpuncak lebih tajam. Masing-masing ekor tersebut tampak mewakili sekitar 2\% hingga 5\% dari luas di bawah distribusi ini.]{0.72}{tDistDF20RightTail1Point65}
  \caption{Kiri: Distribusi $t$ dengan 20 derajat
      kebebasan, dengan daerah di atas 1.65 diarsir.
      Kanan:~Distribusi $t$ dengan 2 derajat kebebasan,
      dengan daerah yang berjarak lebih dari 3 satuan dari 0 diarsir.}
  \label{tDistDF20RightTail1Point65}
\end{figure}

\begin{examplewrap}
\begin{nexample}{Distribusi $t$ dengan 2 derajat kebebasan
    ditampilkan pada panel kanan
    Gambar~\ref{tDistDF20RightTail1Point65}.
    Estimasikan proporsi distribusi yang berjarak lebih
    dari 3~satuan dari rata-rata (di atas atau di bawah).}
  Dengan derajat kebebasan yang begitu sedikit, distribusi $t$
  akan menghasilkan nilai yang jauh berbeda dari hasil distribusi
  normal.
  Pada distribusi normal, luasnya akan sekitar
  0.003 berdasarkan aturan 68-95-99.7.
  Untuk distribusi $t$ dengan $df = 2$, luas pada kedua
  ekor di luar 3~satuan berjumlah 0.0955.
  Luas ini sangat berbeda dari luas yang
  kita peroleh dari distribusi normal.
\end{nexample}
\end{examplewrap}

\begin{exercisewrap}
\begin{nexercise}
Berapa proporsi distribusi $t$ dengan 19 derajat
kebebasan yang berada di atas -1.79 satuan?
Gunakan metode pilihan Anda untuk mencari luas ekor.\footnotemark{}
\end{nexercise}
\end{exercisewrap}
\footnotetext{Kita ingin mencari daerah yang diarsir di \emph{atas}
  -1.79 (pembuatan gambarnya kami serahkan kepada Anda).
  Luas ekor bawah adalah 0.0447,
  sehingga luas daerah atas adalah $1 - 0.0447 = 0.9553$.}

\index{distribusi!t@$t$|)}
\index{distribusi t@distribusi $t$|)}


%\subsection{Conditions for using the $\mathbf{t}$-distribution
%    for inference on a sample mean}
%\label{tDistSolutionToSEProblem}
%
%\noindent%
%To proceed with the $t$-distribution for inference about a single mean, we first check two conditions.
%\begin{description}
%\item[Independence.]
%    We verify this condition just as we did before.
%    We collect a simple random sample, or if the data are from
%    an experiment or random process, we check to the best of our
%    abilities that the observations were independent.
%\item[Sample size.]
%    We use the earlier rule of thumb to evaluate this condition:
%
%    If the sample size $n$ is less than 30
%    and there are no clear outliers in the data,
%    then the sample size condition is satisfied.
%
%    If the sample size $n$ is at least 30
%    and there are no \emph{particularly extreme} outliers,
%    then the sample size condition is satisfied.
%\end{description}
%When examining a sample mean and estimated standard error
%from a sample of $n$ independent and nearly normal observations,
%we use a $t$-distribution with $n - 1$ degrees of freedom~($df$).
%For example, if the sample size was 19, then we would use the
%$t$-distribution with $df = 19 - 1 = 18$ degrees of freedom
%and proceed in a way similar to how we worked with proportions.


\D{\newpage}

\subsection[Interval kepercayaan $t$ satu sampel]
    {Interval kepercayaan $\pmb{t}$ satu sampel}
\label{oneSampleTConfidenceIntervals}

\index{data!lumba-lumba dan merkuri|(}

Mari kita mulai mengenal penerapan distribusi $t$
dalam konteks contoh tentang kandungan merkuri
di otot lumba-lumba.
%Dolphins are at the top of the oceanic food chain, which causes dangerous substances such as mercury to concentrate in their organs and muscles.
Konsentrasi merkuri yang tinggi merupakan masalah penting
baik bagi lumba-lumba
maupun hewan lain, seperti manusia, yang sesekali memakannya.
\captionsetup{width=86mm}

\begin{figure}[h]
\centering
\Figures[Seekor lumba-lumba Risso tampak muncul ke permukaan air. Bagian di depan wajahnya sebagian besar berwarna putih, kemudian tubuhnya bercorak abu-abu dan putih.]{0.8}{rissosDolphin}{rissosDolphin.jpg}  \\
\addvspace{2mm}
\begin{minipage}{\textwidth}
   \caption[rissosDolphinPic]{Seekor lumba-lumba Risso.\vspace{-1mm} \\
   -----------------------------\vspace{-2mm}\\
   {\footnotesize Foto oleh Mike Baird (\oiRedirect{textbook-bairdphotos_com}{www.bairdphotos.com}). \oiRedirect{textbook-CC_BY_2}{lisensi~CC~BY~2.0}.}\vspace{-8mm}}
   \label{rissosDolphin}
\end{minipage}
\stdvspace{}
\end{figure}
\captionsetup{width=\mycaptionwidth}

Kita akan menentukan interval kepercayaan bagi rata-rata kandungan merkuri dalam otot lumba-lumba dengan menggunakan sampel 19 lumba-lumba Risso dari wilayah Taiji di Jepang. Data diringkas dalam Gambar~\ref{summaryStatsOfHgInMuscleOfRissosDolphins}. Nilai minimum dan maksimum yang teramati dapat digunakan untuk menilai apakah terdapat pencilan yang jelas.

\begin{figure}[h]
\centering
\begin{tabular}{ccc cc}
\hline
$n$ & $\bar{x}$ & $s$ & minimum & maksimum \\
19   & 4.4	  & 2.3  & 1.7	       & 9.2 \\
\hline
\end{tabular}
\caption{Ringkasan kandungan merkuri dalam otot
    19 lumba-lumba Risso dari wilayah Taiji.
    Pengukuran dinyatakan dalam mikrogram merkuri per gram basah
    otot ($\mu$g/g basah).}
\label{summaryStatsOfHgInMuscleOfRissosDolphins}
\end{figure}

\begin{examplewrap}
\begin{nexample}{Apakah syarat independensi dan
    normalitas terpenuhi untuk kumpulan~data ini?}
  Observasi-observasi tersebut merupakan sampel acak sederhana,
  sehingga independensi masuk akal.
  Statistik ringkasan dalam
  Gambar~\ref{summaryStatsOfHgInMuscleOfRissosDolphins}
  tidak menunjukkan pencilan yang jelas karena
  semua observasi berada dalam 2.5 simpangan baku
  dari rata-rata.
  Berdasarkan bukti ini, syarat normalitas
  tampaknya masuk akal.
\end{nexample}
\end{examplewrap}

Dalam model normal, kita menggunakan $z^{\star}$ dan galat baku untuk menentukan lebar interval kepercayaan. Kita sedikit menyesuaikan rumus interval kepercayaan ketika menggunakan distribusi $t$:
\begin{align*}
&\text{estimasi titik} \ \pm\  t^{\star}_{df} \times SE
&&\to
&&\bar{x} \ \pm\  t^{\star}_{df} \times \frac{s}{\sqrt{n}}
\end{align*}
%The sample mean is the point estimate of interest.
%The standard error is computed using $SE = s/\sqrt{n}$.

\begin{examplewrap}
\begin{nexample}{Dengan menggunakan statistik ringkasan dalam
    Gambar~\ref{summaryStatsOfHgInMuscleOfRissosDolphins},
    hitung galat baku bagi rata-rata kandungan merkuri
    pada $n = 19$ lumba-lumba tersebut.}
  Kita memasukkan $s$ dan $n$ ke dalam rumus:
  $
  %\begin{align*}
  SE
    = s / \sqrt{n}
    = 2.3 / \sqrt{19}
    = 0.528
  %\end{align*}
  $.
\end{nexample}
\end{examplewrap}

Nilai $t^{\star}_{df}$ merupakan nilai kritis yang kita peroleh berdasarkan
tingkat kepercayaan dan distribusi $t$ dengan $df$ derajat
kebebasan.
Nilai batas itu dicari dengan cara yang sama seperti pada distribusi
normal: kita mencari $t^{\star}_{df}$ sedemikian rupa sehingga
proporsi distribusi $t$ dengan $df$ derajat
kebebasan yang berada dalam jarak $t^{\star}_{df}$
dari 0 sama dengan tingkat kepercayaan yang dikehendaki.

\begin{examplewrap}
\begin{nexample}{Ketika $n = 19$, berapakah banyak
    derajat kebebasan yang sesuai?
    Cari $t^{\star}_{df}$ untuk banyaknya derajat kebebasan ini
    dan tingkat kepercayaan 95\%}
  Banyaknya derajat kebebasan mudah dihitung:
  $df = n - 1 = 18$.
  
  Dengan perangkat lunak statistik, kita mencari nilai kritis yang membuat
  luas ekor atas sama dengan 2.5\%:
  $t^{\star}_{18} = 2.10$.
  Luas di bawah -2.10 juga akan sama dengan 2.5\%.
  Artinya, 95\% distribusi $t$ dengan $df = 18$
  terletak dalam jarak 2.10 satuan dari~0.
\end{nexample}
\end{examplewrap}

%\begin{onebox}{Degrees of freedom for a single sample}
%If the sample has $n$ observations and we are examining a single mean, then we use the $t$-distribution with $df=n-1$ degrees of freedom.
%\end{onebox}

%In our current example, we should use the $t$-distribution
%with $df=19-1=18$ degrees of freedom.
%We can generally identify $t_{18}^{\star}$
%using statistical software.
%Alternatively, we could use the $t$-table in
%Appendix~\ref{tDistributionTable}.
%Generally the value of $t^{\star}_{df}$ is slightly larger
%than what we would get under the normal model with~$z^{\star}$.

\begin{examplewrap}
\begin{nexample}{Hitung dan tafsirkan interval kepercayaan 95\%
    bagi rata-rata kandungan merkuri pada lumba-lumba Risso.}
  Kita dapat menyusun interval kepercayaan sebagai
  \begin{align*}
  \bar{x} \ \pm\  t^{\star}_{18} \times SE
    \quad \to \quad 4.4 \ \pm\  2.10 \times 0.528
    \quad \to \quad (3.29, 5.51)
  \end{align*}
  Kita yakin 95\% bahwa rata-rata kandungan merkuri dalam otot
  lumba-lumba Risso berada antara 3.29 dan 5.51 $\mu$g/gram basah,
  yang dianggap sangat tinggi.
\end{nexample}
\end{examplewrap}

\index{data!lumba-lumba dan merkuri|)}

\begin{onebox}{Menentukan
    interval kepercayaan $\pmb{\MakeLowercase{t}}$
    bagi rata-rata}
  Berdasarkan sampel yang terdiri atas $n$ observasi saling independen
  dan mendekati normal, interval kepercayaan bagi rata-rata
  populasi adalah
  \begin{align*}
  &\text{estimasi titik} \ \pm\  t^{\star}_{df} \times SE
  &&\to
  &&\bar{x} \ \pm\  t^{\star}_{df} \times \frac{s}{\sqrt{n}}
  \end{align*}
  dengan $\bar{x}$ sebagai rata-rata sampel, $t^{\star}_{df}$
  sesuai dengan tingkat kepercayaan dan derajat kebebasan
  $df$, dan $SE$ sebagai galat baku yang diestimasi dari
  sampel.
\end{onebox}

\begin{exercisewrap}
\begin{nexercise} \label{croakerWhiteFishPacificExerConditions}
\index{data!ikan croaker putih dan merkuri|(}
Laman FDA menyediakan sejumlah data tentang kandungan merkuri pada ikan.
Berdasarkan sampel 15 ikan croaker putih (Pasifik),
rata-rata sampel dan simpangan baku masing-masing dihitung sebesar 0.287
dan 0.069 ppm (bagian per sejuta).
Ke-15 observasi tersebut berkisar dari 0.18 hingga 0.41 ppm.
Kita akan menganggap observasi-observasi ini saling independen.
Berdasarkan statistik ringkasan data,
adakah alasan untuk meragukan syarat normalitas
bagi observasi individual?\footnotemark{}
\end{nexercise}
\end{exercisewrap}
\footnotetext{Ukuran sampel kurang dari 30,
  jadi kita memeriksa pencilan yang jelas:
  karena semua observasi berada dalam 2 simpangan baku
  dari rata-rata, tidak ada pencilan yang jelas semacam itu.}

\begin{examplewrap}
\begin{nexample}{Estimasikan galat baku bagi
    $\bar{x} = 0.287$ ppm dengan menggunakan ringkasan data dalam
    Latihan Terarah~\ref{croakerWhiteFishPacificExerConditions}.
    Jika kita hendak menggunakan distribusi $t$ untuk menyusun
    interval kepercayaan 90\% bagi rata-rata sebenarnya dari
    kandungan merkuri, tentukan derajat kebebasan
    dan $t^{\star}_{df}$.}
  \label{croakerWhiteFishPacificExerSEDFTStar}%
  Galat baku: $SE = \frac{0.069}{\sqrt{15}} = 0.0178$.

  Derajat kebebasan: $df = n - 1 = 14$.

  Karena tujuannya adalah interval kepercayaan 90\%,
  kita memilih $t_{14}^{\star}$ sehingga luas kedua ekor
  adalah 0.1:
  $t^{\star}_{14} = 1.76$.
\end{nexample}
\end{examplewrap}

\begin{onebox}{Interval kepercayaan bagi satu rata-rata}
  Setelah menentukan bahwa interval kepercayaan bagi satu rata-rata
  akan berguna dalam suatu penerapan,
  terdapat empat langkah untuk menyusun interval tersebut:
  \begin{description}
  \item[Persiapkan.]
      Tentukan $\bar{x}$, $s$, $n$, dan pilih
      tingkat kepercayaan yang ingin Anda gunakan.
  \item[Periksa.]
      Verifikasi syarat-syarat untuk memastikan bahwa $\bar{x}$
      mendekati normal.
  \item[Hitung.]
      Jika syarat-syarat terpenuhi, hitung $SE$,
      cari $t_{df}^{\star}$, dan susun intervalnya.
  \item[Simpulkan.]
      Tafsirkan interval kepercayaan dalam konteks
      masalah tersebut.
  \end{description}
\end{onebox}

\begin{exercisewrap}
\begin{nexercise}
\label{croakerWhiteFish90ci}
Dengan menggunakan informasi dan hasil Latihan Terarah~\ref{croakerWhiteFishPacificExerConditions} serta Contoh~\ref{croakerWhiteFishPacificExerSEDFTStar}, hitung interval kepercayaan 90\% bagi rata-rata kandungan merkuri pada ikan croaker putih (Pasifik).\footnotemark{}
\end{nexercise}
\end{exercisewrap}
\footnotetext{
  $\bar{x} \ \pm\ t^{\star}_{14} \times SE
      \ \to\  0.287 \ \pm\  1.76 \times 0.0178
      \ \to\ (0.256, 0.318)$.
  Kita yakin 90\% bahwa rata-rata kandungan merkuri
  pada ikan croaker putih (Pasifik) berada antara 0.256 dan 0.318 ppm.}

\begin{exercisewrap}
\begin{nexercise}
Interval kepercayaan 90\% dari
Latihan Terarah~\ref{croakerWhiteFish90ci}
adalah 0.256 ppm hingga 0.318 ppm.
Dapatkah kita mengatakan bahwa 90\% dari ikan croaker putih (Pasifik)
memiliki kadar merkuri antara 0.256 dan 0.318 ppm?\footnotemark{}
\end{nexercise}
\end{exercisewrap}
\footnotetext{
  Tidak. Interval kepercayaan hanya memberikan rentang
  nilai-nilai yang masuk akal bagi suatu parameter populasi,
  dalam hal ini rata-rata populasi.
  Interval itu tidak menjelaskan apa yang mungkin kita amati
  pada observasi individual.}

\index{data!ikan croaker putih dan merkuri|)}

%Now that we've whet \Comment{spelling?} your palette with confidence
%intervals for a mean, let's speed on through to
%hypothesis tests for the mean.


\subsection[Uji $t$ satu sampel]
    {Uji $\pmb{t}$ satu sampel}
\label{oneSampleTTests}
 
\newcommand{\cherryblossomn}{100}
\newcommand{\cherryblossommean}{97.32}
\newcommand{\cherryblossomnull}{93.29}
\newcommand{\cherryblossomsd}{16.98}
\newcommand{\cherryblossomse}{1.70}
\newcommand{\cherryblossomz}{2.37}

\noindent%
Apakah pelari di AS pada umumnya semakin cepat atau semakin lambat dari waktu ke waktu? Kita membahas pertanyaan ini dalam konteks Cherry Blossom Race, yaitu lomba lari 10 mil yang diselenggarakan di Washington, DC setiap~musim semi.

Waktu rata-rata seluruh pelari yang menyelesaikan Cherry Blossom Race pada 2006 adalah \cherryblossomnull{} menit (93 menit dan sekitar 17 detik). Dengan menggunakan data dari \cherryblossomn{} peserta Cherry Blossom Race 2017, kita ingin menentukan apakah para pelari dalam lomba ini semakin cepat atau semakin lambat, dibandingkan dengan kemungkinan lain bahwa tidak terjadi perubahan.

\begin{exercisewrap}
\begin{nexercise}
Hipotesis apa yang sesuai untuk konteks ini?\footnotemark{}
\end{nexercise}
\end{exercisewrap}
\footnotetext{$H_0$: Waktu rata-rata lari 10 mil sama pada 2006 dan 2017. $\mu = \cherryblossomnull{}$ menit. $H_A$: Waktu rata-rata lari 10 mil pada 2017 \emph{berbeda} dari waktu pada 2006. $\mu \neq \cherryblossomnull{}$ menit.}

\begin{exercisewrap}
\begin{nexercise}
Data berasal dari sampel acak sederhana atas seluruh peserta,
sehingga pengamatan-pengamatannya saling independen.
Namun, apakah kita perlu mengkhawatirkan syarat kenormalan?
Lihat Gambar~\ref{run10SampTimeHistogram} yang menampilkan histogram
waktu lari sampel, lalu evaluasilah apakah kita dapat
melanjutkan analisis.\footnotemark{}
\end{nexercise}
\end{exercisewrap}
\footnotetext{Dengan sampel berukuran \cherryblossomn{},
  kita hanya perlu khawatir jika terdapat pencilan yang sangat
  ekstrem.
  Histogram data tidak menunjukkan pencilan yang mengkhawatirkan
  (dan dapat dikatakan tidak menunjukkan pencilan sama sekali).}

\begin{figure}[h]
  \centering
  \Figures[Ditampilkan histogram waktu untuk sampel data Cherry Blossom Race. Data hampir simetris dengan pusat sekitar 100 menit dan simpangan baku kira-kira 15 sampai 20 menit. Semua waktu berada di antara 50 dan 140 menit.]{0.65}{run10SampTimeHistogram}{run17SampTimeHistogram} 
  \caption{Histogram \var{time} untuk sampel data
      Cherry Blossom Race.}
  \label{run10SampTimeHistogram}
\end{figure}

Ketika melakukan uji hipotesis untuk rata-rata satu sampel,
prosesnya hampir sama dengan melakukan uji hipotesis
untuk satu proporsi.
Pertama, kita mencari skor-Z menggunakan nilai teramati,
nilai nol, dan galat baku;
namun, kita menyebutnya \term{skor-T} karena kita menggunakan
distribusi-$t$ untuk menghitung luas ekor.
Kemudian kita mencari nilai-p dengan gagasan yang sama seperti
sebelumnya: cari luas satu ekor di bawah distribusi
sampling, lalu kalikan dua.

\D{\newpage}

%\begin{exampleewrap}
%\begin{nexample}{With independence satisfied and normality
%    not a concern, we can proceed with performing a hypothesis
%    test using the $t$-distribution.
%    The sample mean and sample standard deviation of the
%    sample of \cherryblossomn{} runners from the 2017 Cherry
%    Blossom Race are \cherryblossommean{} and
%    \cherryblossomsd{} minutes, respectively.
%    Recall that the sample size is 100.
%    What is the p-value for the test, and what is your
%    conclusion?}
%\end{nexercise}
%\end{exercisewrap}

\begin{examplewrap}
\begin{nexample}{Karena syarat independensi
    dan kenormalan telah terpenuhi,
    kita dapat melanjutkan uji hipotesis dengan
    distribusi-$t$.
    Rata-rata sampel dan simpangan baku sampel
    dari sampel
    \cherryblossomn{} pelari dalam
    Cherry Blossom Race 2017
    adalah \cherryblossommean{} dan \cherryblossomsd{} menit,
    secara berurutan.
    Ingat bahwa ukuran sampelnya adalah 100
    dan waktu lari rata-rata pada 2006 adalah
    \cherryblossomnull{} menit.
    Carilah statistik uji dan nilai-p.
    Apa kesimpulan Anda?}

  Untuk mencari statistik uji (skor-T),
  terlebih dahulu kita menentukan galat baku:
  \begin{align*}
  SE
    = \cherryblossomsd{} / \sqrt{\cherryblossomn{}}
    = \cherryblossomse{}
  \end{align*}
  Sekarang kita dapat menghitung \emph{skor-T}
  menggunakan rata-rata sampel (\cherryblossommean{}),
  nilai nol (\cherryblossomnull{}), dan $SE$:
  \begin{align*}
  T
    = \frac{\cherryblossommean{} - \cherryblossomnull{}}
        {\cherryblossomse{}}
    = \cherryblossomz{}
  \end{align*}
  Untuk $df = \cherryblossomn{} - 1 = 99$,
  dengan perangkat lunak statistik
  (atau tabel-$t$) kita dapat menentukan bahwa luas satu ekornya adalah 0.01,
  yang kita kalikan dua untuk memperoleh nilai-p:~0.02.

  Karena nilai-p lebih kecil daripada 0.05,
  kita menolak hipotesis nol.
  Artinya, data memberikan bukti kuat bahwa waktu lari
  rata-rata dalam Cherry Blossom Race 2017 berbeda
  dari rata-rata 2006.
  Karena nilai teramati berada di atas nilai nol
  dan kita telah menolak hipotesis nol, kita menyimpulkan
  bahwa pelari dalam lomba tersebut rata-rata lebih lambat pada 2017
  daripada pada 2006.
\end{nexample}
\end{examplewrap}

%%\begin{onebox}{When using a $t$-distribution, we use a T-score (same as Z-score)}
%To help us remember to use the $t$-distribution,
%we use a $T$ to represent the test statistic,
%and we often call this a \term{T-score}.
%The Z-score and T-score are computed in the exact same way
%and are conceptually identical:
%each represents how many standard errors the observed value
%is from the null value.
%%\end{onebox}

\begin{onebox}{Uji hipotesis untuk satu rata-rata}
  Setelah Anda menentukan bahwa uji hipotesis satu rata-rata merupakan
  prosedur yang tepat, terdapat empat langkah untuk menyelesaikan
  pengujian:
  \begin{description}
  \item[Persiapkan.]
      Identifikasi parameter yang menjadi perhatian,
      tuliskan hipotesis,
      tentukan tingkat signifikansi,
      serta identifikasi $\bar{x}$, $s$, dan $n$.
  \item[Periksa.]
      Verifikasi syarat-syarat agar $\bar{x}$ mendekati distribusi normal.
  \item[Hitung.]
      Jika syarat-syarat terpenuhi, hitung $SE$,
      hitung skor-T, dan tentukan nilai-p.
  \item[Simpulkan.]
      Evaluasi uji hipotesis dengan membandingkan nilai-p
      dengan $\alpha$, lalu berikan kesimpulan dalam konteks
      masalah tersebut.
  \end{description}
\end{onebox}

\CalculatorVideos{interval kepercayaan dan uji hipotesis untuk satu rata-rata}


{\input{ch_inference_for_means/TeX/one-sample_means_with_the_t-distribution.tex}}
